\documentclass{article}

\usepackage{amsmath,amssymb}
\usepackage{xurl}
\usepackage[hidelinks]{hyperref}
\setlength{\emergencystretch}{2em}

\input{command}

\title{The Spectral-Radius Identification in Wolf Theorem 6.5}
\author{}
\date{2026-08-24}

\begin{document}
\maketitle
\gapnote{open-gap}{open}

\section{Source statement}

For a linear map $T:\MN{d}\to\MN{d}$, write $\varrho(T)$ for its spectral
radius. For $X\in\MN{d}$, write $X\succeq0$ when $X$ is positive
semidefinite and $X\succ0$ when it is positive definite. Wolf's
Theorem~6.5 states that every positive linear map
$T:\MN{d}\to\MN{d}$ admits a nonzero positive semidefinite eigenvector at its
spectral radius:
\begin{align}
    T(X)
        &=\varrho(T)X,
    \label{eq:wolf_ch6_spectral_radius_eigenvalue_source_eigenvalue_equation}\\
    X
        &\succeq0,
        \qquad X\neq0.
    \label{eq:wolf_ch6_spectral_radius_eigenvalue_source_positive_eigenvector}
\end{align}
This is Theorem~6.5 of Ref.~\cite{Wolf2012Quantum}. The local source is
\path{Notes/WolfNoteTexSource/ch06_spectral_properties.tex},
lines~743--747.

\section{What is formalized}

Two formal results establish separate parts of the source assertion. First,
every positive linear map $E:\MN{D}\to\MN{D}$ has a nonnegative eigenvalue
with a nonzero positive semidefinite eigenvector:
\begin{align}
    E(\rho)
        &=r\rho,
    \label{eq:wolf_ch6_spectral_radius_eigenvalue_formal_existence_equation}\\
    \rho
        &\succeq0,
        \qquad \rho\neq0,
        \qquad r\geq0.
    \label{eq:wolf_ch6_spectral_radius_eigenvalue_formal_existence_conditions}
\end{align}
This is formalized by
\leanid{exists_posSemidef_eigenvector_general}.\footnote{Formal module:
    \doclink{QICLean/Channel/PerronFrobenius/Existence}. The stronger
    declaration \leanid{exists_posSemidef_eigenvector} assumes that $E$
    annihilates no nonzero positive semidefinite matrix and returns $r>0$.}

Second, if a positive map has a positive eigenvalue with a positive definite
eigenvector, then that eigenvalue equals the spectral radius:
\begin{align}
    E(\rho)=r\rho,
    \quad \rho\succ0,
    \quad r>0
        \quad\Longrightarrow\quad
        \varrho(E)=r.
    \label{eq:wolf_ch6_spectral_radius_eigenvalue_positive_definite_identification}
\end{align}
This result requires neither irreducibility nor complete positivity. It is
formalized by
\leanid{spectralRadius_eq_of_posDef_eigenvector_of_positive}.\footnote{Formal
    module: \doclink{QICLean/Channel/Irreducible/SpectralRadius}.}

The declaration
\leanid{exists_wolfTheorem63_of_irreducible_positive} constructs a positive
definite Perron eigenvector and identifies its eigenvalue with the spectral
radius for every irreducible positive map. The source-shaped variant
\leanid{exists_wolfTheorem63_of_irreducible_positive_of_ne_zero} assumes the
map is nonzero and returns an eigenvalue $r>0$.

\section{Comparison}

The existence result
(\ref{eq:wolf_ch6_spectral_radius_eigenvalue_formal_existence_equation}) does
not identify $r$ with the spectral radius. Such an identification is false
for an arbitrary nonnegative eigenvalue of a positive map.

For example, define the completely positive map $E:\MN{2}\to\MN{2}$ by
\begin{align}
    E(X)
        &=X_{11}\one.
    \label{eq:wolf_ch6_spectral_radius_eigenvalue_counterexample_map}
\end{align}
It satisfies
\begin{align}
    E(\one)
        &=\one,
    \label{eq:wolf_ch6_spectral_radius_eigenvalue_counterexample_unit_eigenvector}\\
    \varrho(E)
        &=1.
    \label{eq:wolf_ch6_spectral_radius_eigenvalue_counterexample_spectral_radius}
\end{align}
However, the nonzero positive semidefinite matrix
\begin{align}
    X
        &=\ketbra{2}{2}
    \label{eq:wolf_ch6_spectral_radius_eigenvalue_counterexample_kernel_vector}
\end{align}
lies in the kernel:
\begin{align}
    E(X)
        &=0,
    \label{eq:wolf_ch6_spectral_radius_eigenvalue_counterexample_zero_eigenvalue}\\
    0
        &<\varrho(E).
    \label{eq:wolf_ch6_spectral_radius_eigenvalue_counterexample_strict_gap}
\end{align}
Thus a general existence theorem may return $r=0$ even when the spectral
radius is positive.

The identification theorem
(\ref{eq:wolf_ch6_spectral_radius_eigenvalue_positive_definite_identification})
requires a positive definite eigenvector. Wolf's Theorem~6.5
\cite{Wolf2012Quantum} must also cover reducible maps, for which a
spectral-radius eigenvector may be singular. The irreducible case is complete;
the unresolved step is the extension to arbitrary reducible positive maps.

\section{Routes to a complete proof}

A proof for general positive maps can proceed by either of the following
routes.
\begin{enumerate}
    \item Use the resolvent. For every
        $\lambda>\varrho(T)$, prove that $(\lambda-T)^{-1}$ preserves
        positivity. Then extract a positive semidefinite eigenvector from the
        pole at $\lambda=\varrho(T)$. This route requires holomorphic
        functional calculus and spectral projections.

    \item Use invariant faces of the positive cone. Decompose the positive
        map into irreducible components, apply the Perron--Frobenius theorem
        to each component, and select a component with maximal spectral
        radius. This route requires an invariant-face decomposition for
        positive maps.
\end{enumerate}

Wolf's Theorem~6.3 for irreducible positive maps
\cite{Wolf2012Quantum} is formalized. Its proof follows the lower-maximizer
argument in Equations~(6.31)--(6.33) of Ref.~\cite{Wolf2012Quantum} and the
positive-unital similarity argument; it does not require a separate
pointwise Collatz--Wielandt theorem. The remaining work is exactly the
reducible-map step described above.

\section{Conclusion}

Wolf's Theorem~6.5 \cite{Wolf2012Quantum} remains an open gap. A nonnegative
eigenvalue with a positive semidefinite eigenvector exists for every positive
map, and the spectral-radius identification is established for irreducible
positive maps. The unresolved case is a reducible positive map whose
spectral-radius eigenvector may be singular.

\bibliographystyle{alpha}
\bibliography{references}

\end{document}
